\section{Claims against another party can still affect Apple}
\label{sec:apple-litigation-risk}

The estimate policy leaves a further question: whose claim can give rise to a
cost for Apple? A company can face exposure through an agreement even when the
claim is against someone else. Apple's annual risk discussion says its
agreements sometimes include indemnification provisions that can subject it to
costs and damages in the event of a claim against an indemnified third party
\citep[Risk Factors, printed p.~13]{apple2024}. The direction matters: the
passage describes a possible obligation of Apple arising from another party's
claim, rather than compensation owed to Apple. A model limited to proceedings
naming Apple as defendant would miss that contractual channel.

The statement supplies neither an indemnified party nor a particular agreement.
It also leaves the covered conduct, trigger, duration and limit unknown. To
connect a third-party claim to a dated obligation, the model would need those
terms and the facts that activate them. The existence of indemnification
language somewhere in Apple's agreements cannot establish that a named case
is covered, that a payment is due, or that a liability has been recognised.
Those facts remain distinct from the management loss assessment examined in
Section~\ref{sec:apple-legal}.

The full annual passage also describes unresolved ordinary-course claims and
investigations, an increase over time in their number and alleged magnitude,
and patent claims concerning cellular-enabled products. Apple reports defending
infringement actions in U.S. and international courts, whose plaintiffs
frequently seek broad injunctions and substantial damages. These are issuer
statements about claims and requested relief. They do not identify a complete
matter population, decide the infringement allegations or quantify the
payments and product restrictions that will result.

Defense is a separate source of cost. The filing states, ``Regardless of the
merit of particular claims, defending against litigation or responding to
government investigations can be expensive, time-consuming and disruptive to
the Company's operations.'' Thus a scenario in which Apple successfully
resists a claim need not be a scenario with no defense effort or disruption.
The model therefore needs the response services, their amounts and dates,
and related payment obligations. The damages outcome cannot identify those
expenses or payments.

Settlement can alter recurring operations as well as resolve a dispute. Apple
says it may enter agreements or other arrangements to settle litigation,
without assurance of acceptable terms or that litigation will not occur.
Those agreements can increase cost of sales and operating expenses, change
business practices and limit products and services. A conditional settlement
therefore needs its payment terms and its continuing effects on royalties,
costs or offerings. An adverse resolution can additionally involve monetary
damages, disgorgement, corporate measures or an injunction. The annual passage
explicitly includes a resolution against an indemnified third party above
management's expectations. Neither that possibility nor the separately reported
quarterly settlements supplies a contract, amount or probability for a forecast.

\subsection{Insurance describes a possible recovery, not a net loss}

Apple says it maintains insurance for certain types of claims, while warning
that coverage may be insufficient for all losses or all claim types. The
qualification leaves both the covered population and the amount recoverable
unresolved. A model cannot deduct an assumed insurance payment from damages
or an accrued liability merely because insurance exists. It needs the policy
period, insured party and beneficiary, covered event, exclusions, deductible,
limit, insurer and recovery terms, together with facts about the claim. The
applicable recognition and presentation treatment must also be established
before a recovery amount becomes an accounting entry. The passages identify
none of those amounts or terms.

Other annual insurance statements concern different exposures. The
business-interruption discussion warns of recovery expenditures and lost sales,
especially when critical components have single or limited sources, while
saying insurance for certain types of losses may be insufficient
\citep[Risk Factors, printed p.~6]{apple2024}. The data-security discussion says
coverage is intended for certain aspects of those risks and may not cover all
losses or claim types \citep[Risk Factors, printed p.~12]{apple2024}. The trade
receivable discussion instead says a substantial majority of outstanding trade
receivables lack collateral, third-party bank support or financing arrangements,
or credit insurance \citep[Risk Factors, printed p.~16]{apple2024}. These
statements cannot be combined into one insured share of Apple's exposures.
The receivable qualifier is qualitative; it identifies neither every uncovered
customer nor a percentage. Insurance bought by Apple also does not make these
exposures an insurer's issued-contract measurement problem.

The interim risk preamble incorporates the annual risks and states that there
have been no material changes except those set out below
\citep[Part II, Item 1A, printed p.~19]{apple2025q1}. It preserves the annual
dependency and an express exception; it does not assert identical policies,
claims, coverage or loss amounts at the later date. Keeping the original risk,
the later qualified statement and the separately reported updates allows a
historical query to retain what was actually disclosed. The resulting account
connects direct proceedings, contract indemnities, defense, settlement,
operating restrictions and possible recovery without collapsing them into a
single loss balance. Calculating their consequences still requires the dated
contracts, claim facts and accounting treatment for each channel.
